\chapter*{General Instructions to Students}
\addcontentsline{toc}{chapter}{General Instructions}
\section*{Laboratory Conduct}
\begin{enumerate}[label=\arabic*.]
	\item Students must bring the laboratory record, observation notebook, graph sheets, and a scientific calculator to every session.
	\item The pre-lab assignment for the experiment scheduled for the day must be completed and ready for evaluation before entering the laboratory.
	\item Students shall enter the laboratory only with permission and shall occupy the workbench allotted to them.
	\item Circuit connections must be checked and approved by the faculty-in-charge or laboratory instructor before the power supply is switched on.
	\item Readings must be entered directly into the observation notebook at the time of the experiment; fabricated or estimated readings will not be accepted.
	\item Calculations, graphs, and inferences must be completed and the experiment got corrected on the same day, wherever possible, or within the stipulated time announced by the instructor.
	\item Students must simulate the circuit using the prescribed circuit simulation software prior to hardware implementation, wherever applicable.
	\item Mobile phones shall be kept in silent mode and shall not be used during the conduct of the experiment, except as directed for simulation work.
	\item Any damage to equipment or components due to negligence shall be reported immediately and may be subject to the institution's equipment-damage policy.
	\item Students are responsible for returning all components, probes, and patch cords to their designated locations and for switching off all instruments and the workbench power supply before leaving the laboratory.
\end{enumerate}
\section*{Best Practices for Breadboarding}

Breadboarding is a fundamental skill in analog electronics. A neat and disciplined approach in breadboarding will help achieving the expected results without any hassle. Following these best practices will
help you build reliable circuits and simplify troubleshooting:
\begin{itemize}
	\item \textbf{Plan Your Layout:} Before placing components, mentally (or physically with a drawing)
	plan the layout of your circuit on the breadboard. Group related components.
	\item \textbf{Shortest Connections:} Use short, neat wires for connections. Long, messy wires can
	introduce noise, capacitance, and inductance, and make troubleshooting difficult.
	\item \textbf{Color Coding:} Use different colored wires for different purposes (e.g., red for +Vcc, black
	for ground, blue for signal lines). This greatly aids in tracing connections.
	\item \textbf{Power Rails:} Always connect the power supply rails (+Vcc and GND) to the dedicated
	power buses on the breadboard.
	\item \textbf{Component Orientation:} Pay close attention to the polarity of components like diodes,
	LEDs, and electrolytic capacitors.
	\item \textbf{Secure Connections}: Ensure all component leads and wires are firmly inserted into the
	breadboard holes. Loose connections are a common cause of circuit malfunction.
	\item \textbf{Start Simple:} When troubleshooting, simplify your circuit. Isolate sections and test them
	individually.
	\item \textbf{Component Identification}: Familiarize yourself with common component packages and
	their pinouts (e.g., op-amp pin configurations, transistor pinouts).
	\item \textbf{Avoid Overcrowding:} Don't try to cram too many components into a small area. Use a
	larger breadboard if necessary.
	\item \textbf{Check Continuity:} Before applying power, use a Digital Multimeter (DMM) to check for
	short circuits between power rails and to verify critical connections.
\end{itemize}
\section*{Observations and Data Recording}
Accurate and detailed observations are crucial for a successful lab experience and for writing
comprehensive lab reports.
\begin{itemize}
	\item \textbf{Lab Notebook:} Maintain a dedicated lab notebook.
	\item \textbf{Date and Time:} Always start each entry with the date and time.
	\item \textbf{Experiment Title:} Clearly state the experiment title and number.
	\item \textbf{Circuit Diagrams:} Draw clear and labeled circuit diagrams.
	\item \textbf{Component Values:} Record all component values used in your circuit.
	\item \textbf{Measurements:} Tabulate all measured data with appropriate units.
	\item \textbf{Waveforms:} Sketch or capture screenshots of oscilloscope waveforms, labeling axes,
	amplitudes, and frequencies.
	\item \textbf{Unexpected Results:} Note down any unexpected observations or deviations from
	expected results. This is often where valuable learning occurs.
	\item \textbf{Troubleshooting Steps:} Document any troubleshooting steps taken and their outcomes.
	\item \textbf{Inferences and Discussions:} Write down your inferences and discuss the implications of
	your observations.
\end{itemize}
\section*{Lab Report (Record) Submission}
\begin{enumerate}[label=\arabic*.]
	\item The first thing to note is that the lab report/record is \textbf{NOT A COPY} of the lab manual. The laboratory record is, as its name indicate, a record/report of the experiment done in the lab. The lab manual may include information required for aiding the writing of the report. 
	\item The lab record must contain, for each experiment: aim, components/equipment list, design, circuit diagram (with all component labels/values clearly shown), experimental procedure, observation tables, model graphs, sample calculations, result, and inferences. Please note that it is not required to write the theory and informational notes that appear in this lab manual. 
	\item It is suggested to draw the circuit diagrams and observation tables on the left side of the record, and preferably \textbf{not on a graph page}. Graph pages are for drawing graphs, waveforms and sample/model curves/graphs.
	\item Records must be certified by the faculty-in-charge within one week of completion of the experiment.
	\item Submission of a duly certified record is mandatory for eligibility to appear for the end-semester examination.
\end{enumerate}

\section*{Evaluation of Pre-Lab Assignments}
Two pre-lab assignments are prescribed by the syllabus prior to the regular experiments:
\begin{enumerate}[label=\arabic*.]
	\item Measurement of current, voltage, frequency, and phase shift of a signal in an RC network using a cathode-ray oscilloscope (CRO) / digital storage oscilloscope (DSO).
	\item Introduction to circuit simulation using a circuit simulation software package (e.g., LTspice, PSpice, Multisim, or NgSpice/QUCS).
\end{enumerate}
Students are expected to complete these foundational exercises before attempting the numbered experiments listed in Table~\ref{tab:experiment-list}.